\section{RNN and LSTM --- Concepts and Worked Examples}

% ============================================================
\subsection{Is $h_t$ the Top Hidden Layer?}

Yes. In a stacked RNN, $h_t$ refers to the hidden state of the top hidden layer
(the one closest to the output layer) at time step $t$.

% ============================================================
\subsection{Is $c_t$ a Layer?}

No. $c_t$ is not a layer --- it is an internal memory vector that lives \emph{inside}
the LSTM cell alongside $h_t$. It never connects to the output layer and is never
passed to a higher layer. The distinction is:

\begin{itemize}
    \item $h_t$ is a layer output: passed \textbf{forward} to the next time step and
          \textbf{upward} to the next layer in a stacked architecture.
    \item $c_t$ is an internal state: passed \textbf{forward only} (to the next time step),
          never upward to a higher layer.
\end{itemize}

% ============================================================
\subsection{What Does ``Passed Forward'' and ``Passed Upward'' Mean?}

\begin{itemize}
    \item \textbf{Forward} means passed to the next time step ($t \to t+1$).
          When the RNN finishes processing ``the'', its hidden state is passed forward
          to the step that processes ``hair''.
    \item \textbf{Upward} means passed to the next layer above (in a stacked/deep model).
          In a two-hidden-layer RNN, layer 1's $h_t$ becomes the input to layer 2 at
          the same time step.
\end{itemize}

$c_t$ only goes forward (step to step, staying within the same layer). It never goes
upward to a higher layer.

% ============================================================
\subsection{Do Multiple Hidden Layers in a Stacked LSTM Each Have Their Own $c_t$?}

Yes, and this is by design. In a stacked LSTM with $L$ hidden layers, each layer
maintains its own cell state:

\begin{itemize}
    \item Layer 1 has $c_t^{(1)}$, which passes forward within layer 1 only.
    \item Layer 2 has $c_t^{(2)}$, which passes forward within layer 2 only.
\end{itemize}

The two cell states capture \emph{different levels of abstraction}, not conflicting
memories. Layer 1's $c_t$ tracks low-level patterns (e.g.\ word categories);
layer 2's $c_t$ builds on layer 1's hidden state output and tracks higher-level
patterns (e.g.\ overall sentiment direction). Each hidden layer's $h_t$ is what
travels upward between layers --- not $c_t$.

% ============================================================
\subsection{What Does the Colah LSTM Diagram Show?}

Each green box in the Colah blog diagram represents \textbf{one time step}, not a
stack of layers. So for the sentence ``the hair dryer is great'', the three visible
boxes correspond to three consecutive time steps: processing $x_{t-1}$, $x_t$, and
$x_{t+1}$ respectively.

The two horizontal arrows passing \emph{through} the boxes are $h_t$ and $c_t$
travelling forward from step to step. If you wanted a stacked LSTM (e.g.\ two hidden
layers), that would be drawn as a second row of boxes stacked above these --- not
layers packed inside one box.

The architecture shown in the diagram is therefore: 1 input layer, 1 hidden layer,
1 output layer (a single-layer LSTM unrolled across time steps).

% ============================================================
\subsection{POS Tagging with LSTM: Is Output Produced at Every Step?}

Yes. POS tagging is a many-to-many task, so the LSTM produces one output at every
time step --- one POS tag per word. For ``the hair dryer is great'', the model
produces 5 outputs (one per word), each a probability distribution over all POS tags
via softmax.

% ============================================================
\subsection{What Carries Forward Between Steps?}

Only $h_{t-1}$ and $c_{t-1}$ carry information from the previous step to the current
step. The output layer's softmax probabilities are used only to compute the loss during
training; they do not feed into the next step. The recurrence is:

\begin{align}
    h_t, c_t \quad &\longrightarrow \quad \text{used at step } t+1 \\
    o^{(t)} = \text{softmax}(Vh_t) \quad &\longrightarrow \quad \text{used only for loss, discarded}
\end{align}

% ============================================================
\subsection{Correct POS Tags for the Example Sentence}

For the sentence ``the hair dryer is great'':

\begin{itemize}
    \item the $\to$ DET (determiner)
    \item hair $\to$ NOUN (noun modifier)
    \item dryer $\to$ NOUN
    \item is $\to$ VERB
    \item great $\to$ ADJ
\end{itemize}

% ============================================================
\subsection{Loss Calculation and Backpropagation When a Prediction Is Wrong}

Suppose at step 3 (``dryer'') the model predicts VERB instead of the correct label NOUN.
The cross-entropy loss at that step is:

\begin{equation}
    L^{(3)} = -\log\!\bigl(\hat{o}^{(3)}[\text{NOUN}]\bigr)
\end{equation}

This loss is large because the model assigned low probability to NOUN. However,
the total loss averages across all 5 steps:

\begin{equation}
    L_{\text{total}} = \frac{1}{5}\sum_{t=1}^{5} L^{(t)}
\end{equation}

Backpropagation Through Time (BPTT) then runs \emph{after the full sentence is
processed}. Gradients flow backward through all 5 steps. Because the weight matrices
$W$, $U$, $V$ are \emph{shared} across all steps, their gradient update is the sum of
contributions from every step --- not just the step that made the wrong prediction.

% ============================================================
\subsection{Why Are Weights Updated After the Full Sequence, Not After Each Step?}

Because the weights are shared across all time steps, and because step $t+1$ depends
on step $t$ via the hidden state $h_t$. Specifically:

\begin{itemize}
    \item The loss at step 3 is partly caused by weights that acted at steps 1 and 2
          (since $h_1$ flows into $h_2$, which flows into $h_3$).
    \item You cannot correctly assign blame to those shared weights without seeing
          the total loss across the full sequence.
    \item Updating after each word would compute gradients based on an incomplete
          picture of how the weights contributed to the overall error.
\end{itemize}

This applies equally to vanilla RNNs and LSTMs. The full sequence is treated as one
complete forward pass; one weight update follows at the end.

% ============================================================
\subsection{Comparison with Feedforward Networks and SGD}

In a feedforward network trained with SGD on tabular data (e.g.\ 10 rows, 4 features):

\begin{itemize}
    \item Weights between each layer are \emph{not shared} across rows.
    \item Each row is an independent input with no dependency on other rows.
    \item Therefore: row 1 $\to$ forward pass $\to$ loss $\to$ backprop $\to$ weight
          update. Then row 2, with already-updated weights. And so on.
\end{itemize}

In an RNN, steps within a sequence \emph{are dependent} (via $h_t$), so the
sequence as a whole is treated as one forward pass, with one weight update at the end.

% ============================================================
\subsection{Bidirectional RNN: One Loss or Two?}

One combined loss. The output at each position is produced from both the forward and
backward hidden states fed into a single shared output layer:

\begin{equation}
    o^{(t)} = \text{softmax}\!\bigl(V_f h_f^{(t)} + V_b h_b^{(t)}\bigr)
\end{equation}

This gives one prediction per position $\to$ one loss per position $\to$ one total
loss. Where the two RNNs \emph{separate} is in backpropagation: gradients from the
shared loss flow backward independently through the forward RNN (updating $W_f$,
$U_f$) and the backward RNN (updating $W_b$, $U_b$).

% ============================================================
% ============================================================
\section{Worked Example 1 --- Vanilla RNN POS Tagging}
% ============================================================
% ============================================================

\subsection{Setup}

\textbf{Sentence:} ``the hair dryer is great'' \quad \textbf{Task:} POS tagging
(many-to-many)

\textbf{Architecture:} Input (5 neurons) $\to$ Hidden (5 neurons) $\to$ Output (5 neurons)

\textbf{Vocabulary and one-hot encodings:}
\begin{align*}
    \text{the}   &\to x^{(1)} = [1,0,0,0,0]^\top \\
    \text{hair}  &\to x^{(2)} = [0,1,0,0,0]^\top \\
    \text{dryer} &\to x^{(3)} = [0,0,1,0,0]^\top \\
    \text{is}    &\to x^{(4)} = [0,0,0,1,0]^\top \\
    \text{great} &\to x^{(5)} = [0,0,0,0,1]^\top
\end{align*}

\textbf{POS tag indices:} 0 = DET, 1 = NOUN, 2 = VERB, 3 = ADV, 4 = ADJ

\textbf{True labels:} the $\to$ 0, hair $\to$ 1, dryer $\to$ 1, is $\to$ 2, great $\to$ 4

\subsection{One-Hot True Label Vectors for the Vanilla RNN Example}

The five POS tags are indexed as: 0 = DET, 1 = NOUN, 2 = VERB, 3 = ADV, 4 = ADJ.
The one-hot true label vectors are:

\begin{align*}
    y^{(1)} &= [1, 0, 0, 0, 0]^\top \quad \text{``the''} \to \text{DET (index 0)} \\
    y^{(2)} &= [0, 1, 0, 0, 0]^\top \quad \text{``hair''} \to \text{NOUN (index 1)} \\
    y^{(3)} &= [0, 1, 0, 0, 0]^\top \quad \text{``dryer''} \to \text{NOUN (index 1)} \\
    y^{(4)} &= [0, 0, 1, 0, 0]^\top \quad \text{``is''} \to \text{VERB (index 2)} \\
    y^{(5)} &= [0, 0, 0, 0, 1]^\top \quad \text{``great''} \to \text{ADJ (index 4)}
\end{align*}

Note that ADJ is at index 4, not index 3, because ADV occupies index 3 even though
no word in this sentence is tagged ADV. The output layer still has a neuron for every
POS tag --- the softmax always produces 5 probabilities regardless of which tags
appear in a given sentence.


\subsection{Assumed Weight Matrices}

\begin{equation}
    W = \mathbf{I}_5, \qquad U = 0.5\,\mathbf{I}_5, \qquad V = 0.5\,\mathbf{I}_5, \qquad h^{(0)} = \mathbf{0}
\end{equation}

Since $W = \mathbf{I}_5$, multiplying $W$ by a one-hot vector simply returns that
one-hot vector. Since all matrices are diagonal, $U h^{(t-1)} = 0.5 \cdot h^{(t-1)}$
element-wise, and similarly for $V$.

\subsection{Equations}

\begin{align}
    h^{(t)} &= \tanh\!\bigl(U h^{(t-1)} + W x^{(t)}\bigr) \\
    o^{(t)} &= \text{softmax}\!\bigl(V h^{(t)}\bigr) \\
    L^{(t)} &= -\log\!\bigl(o^{(t)}[\text{true label}]\bigr)
\end{align}

\subsection{Forward Pass}

\textbf{Step 1 --- ``the'':}
\begin{align*}
    U h^{(0)} + W x^{(1)} &= [0,0,0,0,0] + [1,0,0,0,0] = [1,0,0,0,0] \\
    h^{(1)} &= \tanh([1,0,0,0,0]) = [0.7616,0,0,0,0] \\
    V h^{(1)} &= [0.3808,0,0,0,0] \\
    o^{(1)} &= \text{softmax}([0.3808,0,0,0,0]) = [0.2679,0.1830,0.1830,0.1830,0.1830] \\
    L^{(1)} &= -\log(0.2679) = 1.3173 \quad \checkmark \text{ DET predicted correctly}
\end{align*}

\textbf{Step 2 --- ``hair'':}
\begin{align*}
    U h^{(1)} + W x^{(2)} &= [0.3808,0,0,0,0] + [0,1,0,0,0] = [0.3808,1,0,0,0] \\
    h^{(2)} &= \tanh([0.3808,1,0,0,0]) = [0.3634,0.7616,0,0,0] \\
    V h^{(2)} &= [0.1817,0.3808,0,0,0] \\
    o^{(2)} &= \text{softmax}(\cdots) = [0.2118,0.2584,0.1766,0.1766,0.1766] \\
    L^{(2)} &= -\log(0.2584) = 1.3531 \quad \checkmark \text{ NOUN predicted correctly}
\end{align*}

\textbf{Step 3 --- ``dryer'':}
\begin{align*}
    U h^{(2)} + W x^{(3)} &= [0.1817,0.3808,0,0,0] + [0,0,1,0,0] = [0.1817,0.3808,1,0,0] \\
    h^{(3)} &= [0.1797,0.3634,0.7616,0,0] \\
    o^{(3)} &= [0.1900,0.2083,0.2542,0.1737,0.1737] \\
    L^{(3)} &= -\log(0.2083) = 1.5687 \quad \times \text{ VERB predicted, true is NOUN}
\end{align*}

\textbf{Step 4 --- ``is'':}
\begin{align*}
    U h^{(3)} + W x^{(4)} &= [0.0899,0.1817,0.3808,0,0] + [0,0,0,1,0] = [0.0899,0.1817,0.3808,1,0] \\
    h^{(4)} &= [0.0896,0.1797,0.3634,0.7616,0] \\
    o^{(4)} &= [0.1802,0.1885,0.2067,0.2522,0.1723] \\
    L^{(4)} &= -\log(0.2067) = 1.5766 \quad \times \text{ ADV predicted, true is VERB}
\end{align*}

\textbf{Step 5 --- ``great'':}
\begin{align*}
    U h^{(4)} + W x^{(5)} &= [0.0448,0.0899,0.1817,0.3808,0] + [0,0,0,0,1] \\
    h^{(5)} &= [0.0448,0.0896,0.1797,0.3634,0.7616] \\
    o^{(5)} &= [0.1756,0.1795,0.1878,0.2059,0.2512] \\
    L^{(5)} &= -\log(0.2512) = 1.3814 \quad \checkmark \text{ ADJ predicted correctly}
\end{align*}

\textbf{Total Loss:}
\begin{equation}
    L_{\text{total}} = \frac{1.3173 + 1.3531 + 1.5687 + 1.5766 + 1.3814}{5} = 1.4394
\end{equation}

\subsection{Backward Pass (BPTT)}

For cross-entropy loss combined with softmax, the gradient of the loss with respect to
the pre-softmax scores at each step is:
\begin{equation}
    \delta_{\text{out}}^{(t)} = o^{(t)} - y^{(t)}
\end{equation}
where $y^{(t)}$ is the one-hot true label vector.

\textbf{Output layer gradient:}
\begin{align*}
    \delta_{\text{out}}^{(1)} &= [-0.7321,\; 0.1830,\; 0.1830,\; 0.1830,\; 0.1830] \\
    \delta_{\text{out}}^{(2)} &= [\; 0.2118,\;-0.7416,\; 0.1766,\; 0.1766,\; 0.1766] \\
    \delta_{\text{out}}^{(3)} &= [\; 0.1900,\;-0.7917,\; 0.2542,\; 0.1737,\; 0.1737] \\
    \delta_{\text{out}}^{(4)} &= [\; 0.1802,\; 0.1885,\;-0.7933,\; 0.2522,\; 0.1723] \\
    \delta_{\text{out}}^{(5)} &= [\; 0.1756,\; 0.1795,\; 0.1878,\; 0.2059,\;-0.7488]
\end{align*}

\begin{equation}
    \frac{\partial L}{\partial V} = \frac{1}{5}\sum_{t=1}^{5} \delta_{\text{out}}^{(t)} \otimes h^{(t)}
\end{equation}

\textbf{BPTT for $W$ and $U$} (backward from $t=5$ to $t=1$):

At each step, the gradient of the hidden state receives two contributions: one from the
output at that step, and one flowing from the next step. The tanh derivative then
gates it before accumulating into $\nabla W$ and $\nabla U$:
\begin{align*}
    \delta_h^{(t)} &= V^\top \delta_{\text{out}}^{(t)}/5 + \delta_{h,\text{next}} \\
    \delta_{\text{pre}}^{(t)} &= \delta_h^{(t)} \odot \bigl(1 - [h^{(t)}]^2\bigr) \quad
        \leftarrow \text{tanh derivative} \\
    \frac{\partial L}{\partial W} &\mathrel{+}= \delta_{\text{pre}}^{(t)} \otimes x^{(t)} \\
    \frac{\partial L}{\partial U} &\mathrel{+}= \delta_{\text{pre}}^{(t)} \otimes h^{(t-1)} \\
    \delta_{h,\text{next}} &\leftarrow U^\top \delta_{\text{pre}}^{(t)}
\end{align*}

\textbf{Total gradients:}
\begin{equation}
    \nabla V = \begin{bmatrix}
    -0.0845 &  0.0557 &  0.0484 &  0.0402 &  0.0267 \\
    -0.0495 & -0.1605 & -0.1004 &  0.0418 &  0.0273 \\
     0.0373 &  0.0202 & -0.0122 & -0.1072 &  0.0286 \\
     0.0533 &  0.0523 &  0.0522 &  0.0534 &  0.0314 \\
     0.0433 &  0.0323 &  0.0121 & -0.0282 & -0.1141
    \end{bmatrix}
\end{equation}

\begin{equation}
    \nabla W = \begin{bmatrix}
    -0.0240 &  0.0319 &  0.0312 &  0.0266 &  0.0175 \\
    -0.0033 & -0.0431 & -0.0571 &  0.0269 &  0.0178 \\
     0.0266 &  0.0166 & -0.0021 & -0.0610 &  0.0182 \\
     0.0333 &  0.0299 &  0.0245 &  0.0143 &  0.0179 \\
     0.0317 &  0.0267 &  0.0181 &  0.0015 & -0.0314
    \end{bmatrix}
\end{equation}

\begin{equation}
    \nabla U = \begin{bmatrix}
     0.0420 &  0.0366 &  0.0266 &  0.0133 &  0 \\
    -0.0472 & -0.0305 &  0.0269 &  0.0136 &  0 \\
     0.0025 & -0.0205 & -0.0398 &  0.0138 &  0 \\
     0.0359 &  0.0271 &  0.0174 &  0.0136 &  0 \\
     0.0244 &  0.0087 & -0.0103 & -0.0239 &  0
    \end{bmatrix}
\end{equation}

The last column of $\nabla U$ is all zeros because $h^{(0)} = \mathbf{0}$, so no gradient
flows back through the recurrent connection from step 1 to the initial state.

\subsection{Weight Update ($\eta = 0.1$)}

\begin{equation}
    W_{\text{new}} = \begin{bmatrix}
     1.0024 & -0.0032 & -0.0031 & -0.0027 & -0.0018 \\
     0.0003 &  1.0043 &  0.0057 & -0.0027 & -0.0018 \\
    -0.0027 & -0.0017 &  1.0002 &  0.0061 & -0.0018 \\
    -0.0033 & -0.0030 & -0.0025 &  0.9986 & -0.0018 \\
    -0.0032 & -0.0027 & -0.0018 & -0.0002 &  1.0031
    \end{bmatrix}
\end{equation}

\begin{equation}
    U_{\text{new}} = \begin{bmatrix}
     0.4958 & -0.0037 & -0.0027 & -0.0013 &  0 \\
     0.0047 &  0.5030 & -0.0027 & -0.0014 &  0 \\
    -0.0003 &  0.0021 &  0.5040 & -0.0014 &  0 \\
    -0.0036 & -0.0027 & -0.0017 &  0.4986 &  0 \\
    -0.0024 & -0.0009 &  0.0010 &  0.0024 &  0.5
    \end{bmatrix}
\end{equation}

\begin{equation}
    V_{\text{new}} = \begin{bmatrix}
     0.5084 & -0.0056 & -0.0048 & -0.0040 & -0.0027 \\
     0.0049 &  0.5160 &  0.0100 & -0.0042 & -0.0027 \\
    -0.0037 & -0.0020 &  0.5012 &  0.0107 & -0.0029 \\
    -0.0053 & -0.0052 & -0.0052 &  0.4947 & -0.0031 \\
    -0.0043 & -0.0032 & -0.0012 &  0.0028 &  0.5114
    \end{bmatrix}
\end{equation}

% ============================================================
% ============================================================
\section{Worked Example 2 --- Bidirectional RNN POS Tagging}
% ============================================================
% ============================================================

\subsection{Setup}

Same sentence, same vocabulary, same true labels as Example 1.

\textbf{Architecture:} Two RNNs run in parallel on the same sentence. The forward RNN reads
left to right; the backward RNN reads right to left. At each position $t$, the output
combines both hidden states:
\begin{equation}
    o^{(t)} = \text{softmax}\!\bigl(V_f h_f^{(t)} + V_b h_b^{(t)}\bigr)
\end{equation}

\textbf{Assumed weight matrices:}
\begin{equation}
    W_f = \mathbf{I}_5, \quad U_f = 0.5\,\mathbf{I}_5, \quad V_f = 0.5\,\mathbf{I}_5
\end{equation}
\begin{equation}
    W_b = \mathbf{I}_5, \quad U_b = 0.5\,\mathbf{I}_5, \quad V_b = 0.5\,\mathbf{I}_5
\end{equation}
\begin{equation}
    h_f^{(0)} = \mathbf{0}, \qquad h_b^{(6)} = \mathbf{0} \quad \text{(backward initial state)}
\end{equation}

\subsection{Forward RNN (left to right)}

The forward RNN hidden states are identical to those in Example 1:
\begin{align*}
    h_f^{(1)} &= [0.7616,\; 0,\; 0,\; 0,\; 0] \quad \text{(after ``the'')} \\
    h_f^{(2)} &= [0.3634,\; 0.7616,\; 0,\; 0,\; 0] \quad \text{(after ``hair'')} \\
    h_f^{(3)} &= [0.1797,\; 0.3634,\; 0.7616,\; 0,\; 0] \quad \text{(after ``dryer'')} \\
    h_f^{(4)} &= [0.0896,\; 0.1797,\; 0.3634,\; 0.7616,\; 0] \quad \text{(after ``is'')} \\
    h_f^{(5)} &= [0.0448,\; 0.0896,\; 0.1797,\; 0.3634,\; 0.7616] \quad \text{(after ``great'')}
\end{align*}

\subsection{Backward RNN (right to left)}

The backward RNN starts at the end of the sentence with $h_b^{(6)} = \mathbf{0}$ and
processes words from right to left. The equation at each step is:
\begin{equation}
    h_b^{(t)} = \tanh\!\bigl(U_b h_b^{(t+1)} + W_b x^{(t)}\bigr)
\end{equation}

\textbf{Step: $t=5$ --- ``great'' (first backward step):}
\begin{align*}
    U_b h_b^{(6)} + W_b x^{(5)} &= [0,0,0,0,0] + [0,0,0,0,1] = [0,0,0,0,1] \\
    h_b^{(5)} &= \tanh([0,0,0,0,1]) = [0,\; 0,\; 0,\; 0,\; 0.7616]
\end{align*}

\textbf{Step: $t=4$ --- ``is'' (second backward step):}
\begin{align*}
    U_b h_b^{(5)} + W_b x^{(4)} &= [0,0,0,0,0.3808] + [0,0,0,1,0] = [0,0,0,1,0.3808] \\
    h_b^{(4)} &= [0,\; 0,\; 0,\; 0.7616,\; 0.3634]
\end{align*}

\textbf{Step: $t=3$ --- ``dryer'' (third backward step):}
\begin{align*}
    U_b h_b^{(4)} + W_b x^{(3)} &= [0,0,0,0.3808,0.1817] + [0,0,1,0,0] = [0,0,1,0.3808,0.1817] \\
    h_b^{(3)} &= [0,\; 0,\; 0.7616,\; 0.3634,\; 0.1797]
\end{align*}

\textbf{Step: $t=2$ --- ``hair'' (fourth backward step):}
\begin{align*}
    U_b h_b^{(3)} + W_b x^{(2)} &= [0,0,0.3808,0.1817,0.0899] + [0,1,0,0,0] \\
    &= [0,1,0.3808,0.1817,0.0899] \\
    h_b^{(2)} &= [0,\; 0.7616,\; 0.3634,\; 0.1797,\; 0.0896]
\end{align*}

\textbf{Step: $t=1$ --- ``the'' (fifth backward step):}
\begin{align*}
    U_b h_b^{(2)} + W_b x^{(1)} &= [0,0.3808,0.1817,0.0899,0.0448] + [1,0,0,0,0] \\
    &= [1,0.3808,0.1817,0.0899,0.0448] \\
    h_b^{(1)} &= [0.7616,\; 0.3634,\; 0.1797,\; 0.0896,\; 0.0448]
\end{align*}

\noindent\textbf{Note:} $h_b^{(1)}$ (``the'') encodes information about the entire sentence
read from right to left. $h_b^{(5)}$ (``great'') only saw the word ``great'' with no
future context (since it was the first backward step).

\subsection{Combined Output and Loss}

At each position, $V_f h_f^{(t)} + V_b h_b^{(t)}$ combines left-to-right and right-to-left context.

\textbf{Position 1 (``the''):}
\begin{align*}
    V_f h_f^{(1)} &= [0.3808,\; 0,\; 0,\; 0,\; 0] \\
    V_b h_b^{(1)} &= [0.3808,\; 0.1817,\; 0.0899,\; 0.0448,\; 0.0224] \\
    \text{sum}    &= [0.7616,\; 0.1817,\; 0.0899,\; 0.0448,\; 0.0224] \\
    o^{(1)}       &= [0.3293,\; 0.1844,\; 0.1682,\; 0.1608,\; 0.1572] \\
    L^{(1)}       &= -\log(0.3293) = 1.1107 \quad \checkmark \text{ DET}
\end{align*}

\textbf{Position 2 (``hair''):}
\begin{align*}
    V_f h_f^{(2)} + V_b h_b^{(2)} &= [0.1817,\; 0.7616,\; 0.1817,\; 0.0899,\; 0.0448] \\
    o^{(2)} &= [0.1795,\; 0.3206,\; 0.1795,\; 0.1638,\; 0.1566] \\
    L^{(2)} &= -\log(0.3206) = 1.1375 \quad \checkmark \text{ NOUN}
\end{align*}

\textbf{Position 3 (``dryer''):}
\begin{align*}
    V_f h_f^{(3)} + V_b h_b^{(3)} &= [0.0899,\; 0.1817,\; 0.7616,\; 0.1817,\; 0.0899] \\
    o^{(3)} &= [0.1626,\; 0.1782,\; 0.3183,\; 0.1782,\; 0.1626] \\
    L^{(3)} &= -\log(0.1782) = 1.7246 \quad \times \text{ VERB predicted, true is NOUN}
\end{align*}

\textbf{Position 4 (``is''):}
\begin{align*}
    V_f h_f^{(4)} + V_b h_b^{(4)} &= [0.0448,\; 0.0899,\; 0.1817,\; 0.7616,\; 0.1817] \\
    o^{(4)} &= [0.1566,\; 0.1638,\; 0.1795,\; 0.3206,\; 0.1795] \\
    L^{(4)} &= -\log(0.1795) = 1.7174 \quad \times \text{ ADV predicted, true is VERB}
\end{align*}

\textbf{Position 5 (``great''):}
\begin{align*}
    V_f h_f^{(5)} + V_b h_b^{(5)} &= [0.0224,\; 0.0448,\; 0.0899,\; 0.1817,\; 0.7616] \\
    o^{(5)} &= [0.1572,\; 0.1608,\; 0.1682,\; 0.1844,\; 0.3293] \\
    L^{(5)} &= -\log(0.3293) = 1.1107 \quad \checkmark \text{ ADJ}
\end{align*}

\begin{equation}
    L_{\text{total}} = \frac{1.1107 + 1.1375 + 1.7246 + 1.7174 + 1.1107}{5} = 1.3602
\end{equation}

The BiRNN achieves lower total loss than the vanilla RNN (1.3602 vs 1.4394) on this
example because each word now has access to both past and future context.

\subsection{Backward Pass (BPTT)}

There is one combined loss. Gradients from this loss flow backward independently
through two separate RNNs.

\textbf{Gradient of $V_f$ and $V_b$:}
\begin{equation}
    \frac{\partial L}{\partial V_f} = \frac{1}{5}\sum_{t=1}^{5} \delta_{\text{out}}^{(t)} \otimes h_f^{(t)}, \qquad
    \frac{\partial L}{\partial V_b} = \frac{1}{5}\sum_{t=1}^{5} \delta_{\text{out}}^{(t)} \otimes h_b^{(t)}
\end{equation}

\begin{equation}
    \nabla V_f = \begin{bmatrix}
    -0.0791 &  0.0476 &  0.0418 &  0.0353 &  0.0240 \\
    -0.0465 & -0.1544 & -0.1075 &  0.0366 &  0.0245 \\
     0.0369 &  0.0240 & -0.0051 & -0.1127 &  0.0256 \\
     0.0502 &  0.0527 &  0.0571 &  0.0622 &  0.0281 \\
     0.0384 &  0.0301 &  0.0137 & -0.0214 & -0.1022
    \end{bmatrix}
\end{equation}

\begin{equation}
    \nabla V_b = \begin{bmatrix}
    -0.1022 & -0.0214 &  0.0137 &  0.0301 &  0.0384 \\
     0.0281 & -0.0901 & -0.1679 & -0.0559 & -0.0037 \\
     0.0256 &  0.0396 &  0.0676 & -0.0924 & -0.0178 \\
     0.0245 &  0.0366 &  0.0448 &  0.0706 &  0.0622 \\
     0.0240 &  0.0353 &  0.0418 &  0.0476 & -0.0791
    \end{bmatrix}
\end{equation}

\textbf{BPTT for the forward RNN} ($W_f$, $U_f$): gradients flow from $t=5$ back to
$t=1$, exactly as in Example 1. The $\delta_{\text{out}}^{(t)}$ values are different
because the outputs here are from the combined BiRNN model.

\textbf{BPTT for the backward RNN} ($W_b$, $U_b$): gradients flow in the \emph{opposite
direction} of processing. Since the backward RNN processed $t=5$ first and $t=1$ last,
its BPTT flows from $t=1$ \emph{forward} to $t=5$. At each step $t$, the gradient of
$h_b^{(t)}$ includes:
\begin{itemize}
    \item Direct contribution from the output at position $t$: $V_b^\top \delta_{\text{out}}^{(t)}/5$.
    \item Contribution from $h_b^{(t-1)}$ via $U_b$ (since $h_b^{(t)}$ was used to compute
          $h_b^{(t-1)}$ during the forward pass of the backward RNN).
\end{itemize}

\textbf{Total gradients for forward RNN weights:}
\begin{equation}
    \nabla W_f = \begin{bmatrix}
    -0.0224 &  0.0273 &  0.0270 &  0.0233 &  0.0157 \\
    -0.0022 & -0.0414 & -0.0611 &  0.0236 &  0.0160 \\
     0.0258 &  0.0179 & -0.0001 & -0.0641 &  0.0163 \\
     0.0308 &  0.0295 &  0.0262 &  0.0168 &  0.0160 \\
     0.0281 &  0.0248 &  0.0182 &  0.0039 & -0.0282
    \end{bmatrix}
\end{equation}

\begin{equation}
    \nabla U_f = \begin{bmatrix}
     0.0362 &  0.0319 &  0.0235 &  0.0120 &  0 \\
    -0.0480 & -0.0351 &  0.0237 &  0.0121 &  0 \\
     0.0035 & -0.0205 & -0.0429 &  0.0124 &  0 \\
     0.0365 &  0.0290 &  0.0186 &  0.0122 &  0 \\
     0.0236 &  0.0102 & -0.0073 & -0.0215 &  0
    \end{bmatrix}
\end{equation}

\textbf{Total gradients for backward RNN weights:}
\begin{equation}
    \nabla W_b = \begin{bmatrix}
    -0.0282 &  0.0039 &  0.0182 &  0.0248 &  0.0281 \\
     0.0160 & -0.0252 & -0.0948 & -0.0310 &  0.0006 \\
     0.0163 &  0.0226 &  0.0181 & -0.0730 & -0.0197 \\
     0.0160 &  0.0236 &  0.0257 &  0.0189 &  0.0279 \\
     0.0157 &  0.0233 &  0.0270 &  0.0273 & -0.0224
    \end{bmatrix}
\end{equation}

\begin{equation}
    \nabla U_b = \begin{bmatrix}
     0      & -0.0215 & -0.0073 &  0.0102 &  0.0236 \\
     0      &  0.0122 & -0.0134 & -0.0784 & -0.0611 \\
     0      &  0.0124 &  0.0232 &  0.0250 & -0.0435 \\
     0      &  0.0121 &  0.0237 &  0.0310 &  0.0294 \\
     0      &  0.0120 &  0.0235 &  0.0319 &  0.0362
    \end{bmatrix}
\end{equation}

The first column of $\nabla U_b$ is all zeros because $h_b^{(6)} = \mathbf{0}$ (the
backward RNN's initial state), mirroring how $\nabla U$'s last column was zero in
Example 1 for the same reason.

\subsection{Weight Update ($\eta = 0.1$)}

\begin{equation}
    W_{f,\text{new}} = \begin{bmatrix}
     1.0022 & -0.0027 & -0.0027 & -0.0023 & -0.0016 \\
     0.0002 &  1.0041 &  0.0061 & -0.0024 & -0.0016 \\
    -0.0026 & -0.0018 &  1.0000 &  0.0064 & -0.0016 \\
    -0.0031 & -0.0029 & -0.0026 &  0.9983 & -0.0016 \\
    -0.0028 & -0.0025 & -0.0018 & -0.0004 &  1.0028
    \end{bmatrix}
\end{equation}

\begin{equation}
    U_{f,\text{new}} = \begin{bmatrix}
     0.4964 & -0.0032 & -0.0023 & -0.0012 &  0 \\
     0.0048 &  0.5035 & -0.0024 & -0.0012 &  0 \\
    -0.0004 &  0.0020 &  0.5043 & -0.0012 &  0 \\
    -0.0036 & -0.0029 & -0.0019 &  0.4988 &  0 \\
    -0.0024 & -0.0010 &  0.0007 &  0.0021 &  0.5
    \end{bmatrix}
\end{equation}

\begin{equation}
    W_{b,\text{new}} = \begin{bmatrix}
     1.0028 & -0.0004 & -0.0018 & -0.0025 & -0.0028 \\
    -0.0016 &  1.0025 &  0.0095 &  0.0031 & -0.0001 \\
    -0.0016 & -0.0023 &  0.9982 &  0.0073 &  0.0020 \\
    -0.0016 & -0.0024 & -0.0026 &  0.9981 & -0.0028 \\
    -0.0016 & -0.0023 & -0.0027 & -0.0027 &  1.0022
    \end{bmatrix}
\end{equation}

\begin{equation}
    U_{b,\text{new}} = \begin{bmatrix}
     0.5    &  0.0021 &  0.0007 & -0.0010 & -0.0024 \\
     0      &  0.4988 &  0.0013 &  0.0078 &  0.0061 \\
     0      & -0.0012 &  0.4977 & -0.0025 &  0.0043 \\
     0      & -0.0012 & -0.0024 &  0.4969 & -0.0029 \\
     0      & -0.0012 & -0.0023 & -0.0032 &  0.4964
    \end{bmatrix}
\end{equation}

\begin{equation}
    V_{f,\text{new}} = \begin{bmatrix}
     0.5079 & -0.0048 & -0.0042 & -0.0035 & -0.0024 \\
     0.0046 &  0.5154 &  0.0107 & -0.0037 & -0.0024 \\
    -0.0037 & -0.0024 &  0.5005 &  0.0113 & -0.0026 \\
    -0.0050 & -0.0053 & -0.0057 &  0.4938 & -0.0028 \\
    -0.0038 & -0.0030 & -0.0014 &  0.0021 &  0.5102
    \end{bmatrix}
\end{equation}

\begin{equation}
    V_{b,\text{new}} = \begin{bmatrix}
     0.5102 &  0.0021 & -0.0014 & -0.0030 & -0.0038 \\
    -0.0028 &  0.5090 &  0.0168 &  0.0056 &  0.0004 \\
    -0.0026 & -0.0040 &  0.4932 &  0.0092 &  0.0018 \\
    -0.0024 & -0.0037 & -0.0045 &  0.4929 & -0.0062 \\
    -0.0024 & -0.0035 & -0.0042 & -0.0048 &  0.5079
    \end{bmatrix}
\end{equation}

% ============================================================
% ============================================================
\section{Worked Example 3 --- LSTM POS Tagging}
% ============================================================
% ============================================================

\subsection{Setup}

Same sentence, vocabulary, and true labels as Examples 1 and 2.

\textbf{Architecture:} Input (5) $\to$ LSTM hidden (5) $\to$ Output (5)

\textbf{LSTM Equations} (split-form notation, no bias):
\begin{align}
    f_t &= \sigma(W_{f,h}\, h_{t-1} + W_{f,x}\, x_t) \quad \text{(forget gate)} \\
    g_t &= \tanh(W_{g,h}\, h_{t-1} + W_{g,x}\, x_t) \quad \text{(candidate)} \\
    i_t &= \sigma(W_{i,h}\, h_{t-1} + W_{i,x}\, x_t) \quad \text{(input gate)} \\
    o_t &= \sigma(W_{o,h}\, h_{t-1} + W_{o,x}\, x_t) \quad \text{(output gate)} \\
    c_t &= f_t \odot c_{t-1} + i_t \odot g_t \quad \text{(cell state)} \\
    h_t &= o_t \odot \tanh(c_t) \quad \text{(hidden state)}
\end{align}

\textbf{Note on notation:} These are equivalent to Colah's concatenation form
$W \cdot [h_{t-1}, x_t]$ where $W = [W_h \mid W_x]$ is the concatenated matrix.
Both forms produce identical results.

\textbf{Assumed weight matrices} (all diagonal, no biases):
\begin{equation}
    W_{f,h} = 0.4\,\mathbf{I}_5, \quad W_{f,x} = 0.3\,\mathbf{I}_5 \quad \text{(forget gate)}
\end{equation}
\begin{equation}
    W_{g,h} = 0.2\,\mathbf{I}_5, \quad W_{g,x} = 0.7\,\mathbf{I}_5 \quad \text{(candidate)}
\end{equation}
\begin{equation}
    W_{i,h} = 0.5\,\mathbf{I}_5, \quad W_{i,x} = 0.4\,\mathbf{I}_5 \quad \text{(input gate)}
\end{equation}
\begin{equation}
    W_{o,h} = 0.3\,\mathbf{I}_5, \quad W_{o,x} = 0.4\,\mathbf{I}_5 \quad \text{(output gate)}
\end{equation}
\begin{equation}
    V = 0.5\,\mathbf{I}_5, \qquad h_0 = \mathbf{0}, \qquad c_0 = \mathbf{0}
\end{equation}

Since all matrices are diagonal, multiplying by a vector is equivalent to element-wise
scaling. For example, $0.4\,\mathbf{I}_5 \cdot h_{t-1} = 0.4 \cdot h_{t-1}$ element-wise.

\subsection{Forward Pass}

\textbf{Step 1 --- ``the'':}
\begin{align*}
    h_0 &= [0,0,0,0,0], \quad c_0 = [0,0,0,0,0], \quad x^{(1)} = [1,0,0,0,0] \\[4pt]
    f_1 &= \sigma(0.4 \cdot [0,0,0,0,0] + 0.3 \cdot [1,0,0,0,0]) \\
        &= \sigma([0.3,0,0,0,0]) = [0.5744,\;0.5000,\;0.5000,\;0.5000,\;0.5000] \\[4pt]
    g_1 &= \tanh(0.2 \cdot [0,0,0,0,0] + 0.7 \cdot [1,0,0,0,0]) \\
        &= \tanh([0.7,0,0,0,0]) = [0.6044,\;0,\;0,\;0,\;0] \\[4pt]
    i_1 &= \sigma(0.5 \cdot [0,0,0,0,0] + 0.4 \cdot [1,0,0,0,0]) \\
        &= \sigma([0.4,0,0,0,0]) = [0.5987,\;0.5000,\;0.5000,\;0.5000,\;0.5000] \\[4pt]
    o_1 &= \sigma(0.3 \cdot [0,0,0,0,0] + 0.4 \cdot [1,0,0,0,0]) \\
        &= \sigma([0.4,0,0,0,0]) = [0.5987,\;0.5000,\;0.5000,\;0.5000,\;0.5000] \\[4pt]
    c_1 &= f_1 \odot c_0 + i_1 \odot g_1
         = [0,0,0,0,0] + [0.5987 \times 0.6044,\; 0.5\times 0,\;\ldots] \\
        &= [0.3618,\;0,\;0,\;0,\;0] \\[4pt]
    h_1 &= o_1 \odot \tanh(c_1)
         = [0.5987,\;0.5,\;0.5,\;0.5,\;0.5] \odot [0.3468,\;0,\;0,\;0,\;0] \\
        &= [0.2076,\;0,\;0,\;0,\;0] \\[4pt]
    o^{(1)} &= \text{softmax}(V h_1) = \text{softmax}([0.1038,0,0,0,0]) \\
            &= [0.2171,\;0.1957,\;0.1957,\;0.1957,\;0.1957] \\
    L^{(1)} &= -\log(0.2171) = 1.5273 \quad \checkmark \text{ DET predicted correctly}
\end{align*}

\textbf{Step 2 --- ``hair'':}
\begin{align*}
    h_1 &= [0.2076,0,0,0,0], \quad c_1 = [0.3618,0,0,0,0] \\
    &\quad x^{(2)} = [0,1,0,0,0] \\[4pt]
    f_2 &= \sigma(0.4 \cdot h_1 + 0.3 \cdot x^{(2)}) \\
         &= \sigma([0.0831,\;0.3,\;0,\;0,\;0])
         = [0.5208,\;0.5744,\;0.5000,\;0.5000,\;0.5000] \\[4pt]
    g_2 &= \tanh(0.2 \cdot h_1 + 0.7 \cdot x^{(2)}) \\
         &= \tanh([0.0415,\;0.7,\;0,\;0,\;0])
         = [0.0415,\;0.6044,\;0,\;0,\;0] \\[4pt]
    i_2 &= \sigma(0.5 \cdot h_1 + 0.4 \cdot x^{(2)}) \\
         &= \sigma([0.1038,\;0.4,\;0,\;0,\;0])
         = [0.5259,\;0.5987,\;0.5000,\;0.5000,\;0.5000] \\[4pt]
    o_2 &= \sigma(0.3 \cdot h_1 + 0.4 \cdot x^{(2)}) \\
         &= \sigma([0.0623,\;0.4,\;0,\;0,\;0])
         = [0.5156,\;0.5987,\;0.5000,\;0.5000,\;0.5000] \\[4pt]
    c_2 &= f_2 \odot c_1 + i_2 \odot g_2 \\
        &= [0.5208 \times 0.3618,\;0.5744\times 0,\;\ldots]
         + [0.5259 \times 0.0415,\;0.5987 \times 0.6044,\;\ldots] \\
        &= [0.1884,\;0,\;0,\;0,\;0] + [0.0218,\;0.3618,\;0,\;0,\;0] \\
        &= [0.2103,\;0.3618,\;0,\;0,\;0] \\[4pt]
    h_2 &= o_2 \odot \tanh(c_2)
         = [0.5156,\;0.5987,\;\ldots] \odot [0.2072,\;0.3468,\;\ldots] \\
         &= [0.1068,\;0.2076,\;0,\;0,\;0] \\[4pt]
    o^{(2)} &= \text{softmax}([0.0534,\;0.1038,\;0,\;0,\;0]) \\
             &= [0.2043,\;0.2148,\;0.1936,\;0.1936,\;0.1936] \\
    L^{(2)} &= -\log(0.2148) = 1.5379 \\
    &\quad \checkmark \text{ NOUN predicted correctly}
\end{align*}

\textbf{Step 3 --- ``dryer'':}
\begin{align*}
    h_2 &= [0.1068,0.2076,0,0,0], \quad c_2 = [0.2103,0.3618,0,0,0] \\[4pt]
    f_3 &= \sigma([0.0427,\;0.0831,\;0.3,\;0,\;0])
         = [0.5107,\;0.5208,\;0.5744,\;0.5,\;0.5] \\
    g_3 &= \tanh([0.0214,\;0.0415,\;0.7,\;0,\;0])
         = [0.0214,\;0.0415,\;0.6044,\;0,\;0] \\
    i_3 &= \sigma([0.0534,\;0.1038,\;0.4,\;0,\;0])
         = [0.5134,\;0.5259,\;0.5987,\;0.5,\;0.5] \\
    o_3 &= \sigma([0.0320,\;0.0623,\;0.4,\;0,\;0])
         = [0.5080,\;0.5156,\;0.5987,\;0.5,\;0.5] \\[4pt]
    c_3 &= f_3 \odot c_2 + i_3 \odot g_3 \\
        &= [0.1074,\;0.1884,\;0,\;0,\;0] + [0.0110,\;0.0218,\;0.3618,\;0,\;0] \\
        &= [0.1183,\;0.2103,\;0.3618,\;0,\;0] \\[4pt]
    h_3 &= o_3 \odot \tanh(c_3) = [0.0598,\;0.1068,\;0.2076,\;0,\;0] \\[4pt]
    o^{(3)} &= \text{softmax}([0.0299,\;0.0534,\;0.1038,\;0,\;0]) \\
             &= [0.1984,\;0.2031,\;0.2136,\;0.1925,\;0.1925] \\
    L^{(3)} &= -\log(0.2031) = 1.5942 \quad \times \text{ VERB predicted, true is NOUN}
\end{align*}

\textbf{Step 4 --- ``is'':}
\begin{align*}
    h_3 &= [0.0598,0.1068,0.2076,0,0], \quad c_3 = [0.1183,0.2103,0.3618,0,0] \\[4pt]
    f_4 &= \sigma([0.0239,0.0427,0.0831,0.3,0])
         = [0.5060,0.5107,0.5208,0.5744,0.5] \\
    g_4 &= \tanh([0.0120,0.0214,0.0415,0.7,0])
         = [0.0120,0.0214,0.0415,0.6044,0] \\
    i_4 &= \sigma([0.0299,0.0534,0.1038,0.4,0])
         = [0.5075,0.5134,0.5259,0.5987,0.5] \\
    o_4 &= \sigma([0.0180,0.0320,0.0623,0.4,0])
         = [0.5045,0.5080,0.5156,0.5987,0.5] \\[4pt]
    c_4 &= [0.0599,0.1074,0.1884,0,0] + [0.0061,0.0110,0.0218,0.3618,0] \\
        &= [0.0659,0.1183,0.2103,0.3618,0] \\[4pt]
    h_4 &= [0.0332,0.0598,0.1068,0.2076,0] \\[4pt]
    o^{(4)} &= \text{softmax}([0.0166,0.0299,0.0534,0.1038,0]) \\
             &= [0.1951,0.1977,0.2024,0.2129,0.1919] \\
    L^{(4)} &= -\log(0.2024) = 1.5974 \quad \times \text{ ADV predicted, true is VERB}
\end{align*}

\textbf{Step 5 --- ``great'':}
\begin{align*}
    h_4 &= [0.0332,0.0598,0.1068,0.2076,0] \\
    &\quad c_4 = [0.0659,0.1183,0.2103,0.3618,0] \\[4pt]
    f_5 &= \sigma([0.0133,0.0239,0.0427,0.0831,0.3])
         = [0.5033,0.5060,0.5107,0.5208,0.5744] \\
    g_5 &= \tanh([0.0066,0.0120,0.0214,0.0415,0.7])
         = [0.0066,0.0120,0.0214,0.0415,0.6044] \\
    i_5 &= \sigma([0.0166,0.0299,0.0534,0.1038,0.4])
         = [0.5042,0.5075,0.5134,0.5259,0.5987] \\
    o_5 &= \sigma([0.0100,0.0180,0.0320,0.0623,0.4])
         = [0.5025,0.5045,0.5080,0.5156,0.5987] \\[4pt]
    c_5 &= [0.0332,0.0599,0.1074,0.1884,0] + [0.0033,0.0061,0.0110,0.0218,0.3618] \\
        &= [0.0365,0.0659,0.1183,0.2103,0.3618] \\[4pt]
    h_5 &= [0.0184,0.0332,0.0598,0.1068,0.2076] \\[4pt]
    o^{(5)} &= \text{softmax}([0.0092,0.0166,0.0299,0.0534,0.1038]) \\
             &= [0.1933,0.1948,0.1974,0.2021,0.2125] \\
    L^{(5)} &= -\log(0.2125) = 1.5488 \\
    &\quad \checkmark \text{ ADJ predicted correctly}
\end{align*}

\begin{equation}
    L_{\text{total}} = \frac{1.5273 + 1.5379 + 1.5942 + 1.5974 + 1.5488}{5} = 1.5611
\end{equation}

\subsection{Backward Pass (BPTT)}

\textbf{Output layer gradient:}
\begin{align*}
    \delta_{\text{out}}^{(1)} &= [-0.7829,\; 0.1957,\; 0.1957,\; 0.1957,\; 0.1957] \\
    \delta_{\text{out}}^{(2)} &= [\; 0.2043,\;-0.7852,\; 0.1936,\; 0.1936,\; 0.1936] \\
    \delta_{\text{out}}^{(3)} &= [\; 0.1984,\;-0.7969,\; 0.2136,\; 0.1925,\; 0.1925] \\
    \delta_{\text{out}}^{(4)} &= [\; 0.1951,\; 0.1977,\;-0.7976,\; 0.2129,\; 0.1919] \\
    \delta_{\text{out}}^{(5)} &= [\; 0.1933,\; 0.1948,\; 0.1974,\; 0.2021,\;-0.7875]
\end{align*}

\textbf{LSTM BPTT chain} (from $t=5$ back to $t=1$):

At each step $t$, the gradient flows through the output gate and the cell state. The
cell state gradient is the LSTM's gradient highway:
\begin{align}
    \delta_h^{(t)} &= V^\top \delta_{\text{out}}^{(t)}/T + \delta_{h,\text{next}} \\
    \delta_{o}^{(t)} &= \delta_h^{(t)} \odot \tanh(c_t) \\
    \delta_{o,\text{raw}}^{(t)} &= \delta_o^{(t)} \odot o_t \odot (1 - o_t) \quad \leftarrow \sigma' \\
    \delta_c^{(t)} &= \delta_h^{(t)} \odot o_t \odot (1 - \tanh^2(c_t)) + \delta_{c,\text{next}} \\
    \delta_{f,\text{raw}}^{(t)} &= (\delta_c^{(t)} \odot c_{t-1}) \odot f_t \odot (1 - f_t) \quad \leftarrow \sigma' \\
    \delta_{g,\text{raw}}^{(t)} &= (\delta_c^{(t)} \odot i_t) \odot (1 - g_t^2) \quad \leftarrow \tanh' \\
    \delta_{i,\text{raw}}^{(t)} &= (\delta_c^{(t)} \odot g_t) \odot i_t \odot (1 - i_t) \quad \leftarrow \sigma'
\end{align}

Note: the gradient of $c_t$ with respect to $c_{t-1}$ passes only through the forget
gate $f_t$ (no matrix multiplication, no tanh squashing) --- this is the gradient highway
that prevents vanishing gradients.

Key values during BPTT ($t = 5$ down to $t = 1$):

\begin{center}
\resizebox{\textwidth}{!}{%
\begin{tabular}{c|l|l|l|l}
$t$ & $\delta_{o,\text{raw}}^{(t)}$ & $\delta_c^{(t)}$ & $\delta_{f,\text{raw}}^{(t)}$ & $\delta_{g,\text{raw}}^{(t)}$ \\
\hline
5 & $[0.0002,0.0003,0.0006,0.0010,-0.0066]$ & $[0.0097,0.0098,0.0099,0.0100,-0.0415]$ & $[\approx 0]$ & $[0.0049,0.0050,0.0051,0.0052,-0.0158]$ \\
4 & $[0.0003,0.0006,-0.0041,0.0019,0.0000]$ & $[0.0152,0.0155,-0.0336,0.0173,-0.0183]$ & $[0.0005,0.0008,-0.0030,\approx 0]$ & $[0.0077,0.0079,-0.0176,0.0066,-0.0091]$ \\
3 & $[0.0006,-0.0040,0.0013,\approx 0]$ & $[0.0186,-0.0304,-0.0095,0.0212,-0.0004]$ & $[0.0010,-0.0027,\approx 0]$ & $[0.0095,-0.0159,-0.0036,0.0106,-0.0002]$ \\
2 & $[0.0012,-0.0070,\approx 0]$ & $[0.0208,-0.0601,0.0037,0.0213,0.0094]$ & $[0.0019,\approx 0]$ & $[0.0109,-0.0229,0.0019,0.0107,0.0047]$ \\
1 & $[-0.0062,\approx 0]$ & $[-0.0286,-0.0303,0.0118,0.0215,0.0150]$ & $[\approx 0]$ & $[-0.0109,-0.0151,0.0059,0.0108,0.0075]$
\end{tabular}%
}
\end{center}

\subsection{Total Gradients}

\begin{equation}
    \nabla V = \begin{bmatrix}
    -0.0238 &  0.0163 &  0.0147 &  0.0122 &  0.0080 \\
    -0.0162 & -0.0460 & -0.0265 &  0.0124 &  0.0081 \\
     0.0102 &  0.0044 & -0.0058 & -0.0289 &  0.0082 \\
     0.0167 &  0.0160 &  0.0150 &  0.0132 &  0.0084 \\
     0.0130 &  0.0092 &  0.0027 & -0.0089 & -0.0327
    \end{bmatrix}
\end{equation}

\begin{equation}
    \nabla W_{g,x} = \begin{bmatrix}
    -0.0109 &  0.0109 &  0.0095 &  0.0077 &  0.0049 \\
    -0.0151 & -0.0229 & -0.0159 &  0.0079 &  0.0050 \\
     0.0059 &  0.0019 & -0.0036 & -0.0176 &  0.0051 \\
     0.0108 &  0.0107 &  0.0106 &  0.0066 &  0.0052 \\
     0.0075 &  0.0047 & -0.0002 & -0.0091 & -0.0158
    \end{bmatrix}
\end{equation}

The gradients for the gate weight matrices involving $h$ ($W_{f,h}$, $W_{g,h}$,
$W_{i,h}$, $W_{o,h}$) are much smaller in magnitude than the $x$-side gradients
because the hidden states are small in early steps (the model is building up memory).
The $x$-side gradients ($W_{f,x}$, $W_{g,x}$, $W_{i,x}$, $W_{o,x}$) are larger
because the input signal is binary (one-hot) and directly controls what the LSTM writes
into memory at each step.

\subsection{Weight Update ($\eta = 0.1$)}

\begin{equation}
    V_{\text{new}} = \begin{bmatrix}
     0.5024 & -0.0016 & -0.0015 & -0.0012 & -0.0008 \\
     0.0016 &  0.5046 &  0.0027 & -0.0012 & -0.0008 \\
    -0.0010 & -0.0004 &  0.5006 &  0.0029 & -0.0008 \\
    -0.0017 & -0.0016 & -0.0015 &  0.4987 & -0.0008 \\
    -0.0013 & -0.0009 & -0.0003 &  0.0009 &  0.5033
    \end{bmatrix}
\end{equation}

\begin{equation}
    W_{g,x,\text{new}} = \begin{bmatrix}
     0.7011 & -0.0011 & -0.0010 & -0.0008 & -0.0005 \\
     0.0015 &  0.7023 &  0.0016 & -0.0008 & -0.0005 \\
    -0.0006 & -0.0002 &  0.7004 &  0.0018 & -0.0005 \\
    -0.0011 & -0.0011 & -0.0011 &  0.6993 & -0.0005 \\
    -0.0007 & -0.0005 &  0.0000 &  0.0009 &  0.7016
    \end{bmatrix}
\end{equation}

All other updated gate weight matrices are close to their initial values (changes in
the third or fourth decimal place), which reflects that the LSTM requires many more
training iterations to show significant weight movement. The complete updated matrices
for all 8 gate weight pairs are as follows:

\begin{align*}
    W_{f,h,\text{new}} &\approx 0.4\,\mathbf{I}_5 \quad \text{(changes } \leq 0.0006\text{)} \\
    W_{f,x,\text{new}} &\approx 0.3\,\mathbf{I}_5 \quad \text{(changes } \leq 0.0003\text{)} \\
    W_{g,h,\text{new}} &\approx 0.2\,\mathbf{I}_5 \quad \text{(changes } \leq 0.0006\text{)} \\
    W_{i,h,\text{new}} &\approx 0.5\,\mathbf{I}_5 \quad \text{(changes } \leq 0.0018\text{)} \\
    W_{i,x,\text{new}} &\approx 0.4\,\mathbf{I}_5 \quad \text{(changes } \leq 0.0060\text{)} \\
    W_{o,h,\text{new}} &\approx 0.3\,\mathbf{I}_5 \quad \text{(changes } \leq 0.0018\text{)} \\
    W_{o,x,\text{new}} &\approx 0.4\,\mathbf{I}_5 \quad \text{(changes } \leq 0.0070\text{)}
\end{align*}

\subsection{Key Observations from the LSTM Walkthrough}

\begin{itemize}
    \item \textbf{Cell state accumulates gradually.} At step 5, {\scriptsize $c_5 = [0.0365, 0.0659, 0.1183, 0.2103, 0.3618]$}. Each new word's information is added through $i_t \odot g_t$, while old information is preserved through $f_t \odot c_{t-1}$. The diagonal weight structure means each position in the cell state tracks the corresponding word independently.

    \item \textbf{Forget gate near 0.5 throughout.} Because $h_{t-1}$ is small and the initial weights are small, the forget gate values hover near $\sigma(0) = 0.5$ for most entries. This means roughly half of the old cell state survives at each step. In a trained model, the forget gate would learn to output values closer to 0 or 1 depending on context.

    \item \textbf{Hidden state is much smaller than cell state.} The output gate multiplies $\tanh(c_t)$ by approximately 0.5 (since $W_{o,h}$ and $W_{o,x}$ are small), so $h_t \approx 0.5 \cdot \tanh(c_t)$. This demonstrates how the output gate filters what the cell state exposes to the output layer.

    \item \textbf{Gradient highway in action.} During BPTT, $\delta_c^{(t)}$ at step 5 is\\
{\scriptsize $[-0.0286, -0.0303, 0.0118, 0.0215, 0.0150]$} at step 1 --- a non-vanishing gradient that reaches the earliest time step. In the vanilla RNN, the corresponding gradient at step 1 was much closer to zero due to repeated tanh squashing.
\end{itemize}